\section*{अध्याय \ref{ch:b2:phanerozoic-earth} --- फ़ैनेरोज़ोइक पृथ्वी}

\begin{solution}{exo:b2:phanerozoic-earth:1}
(क) पेरिस (\ang{5.1}~N) और न्यूयॉर्क (\ang{3.6}~S); केप टाउन और सिडनी \ang{60}~S के पास थे।
(ख) केप टाउन: $61.9 - 33.9 = \ang{28.0}$, $28.0 \times 111.2 = \qty{3114}{km}$; न्यूयॉर्क:
$40.7 + 3.6 = \ang{44.3}$, \qty{4926}{km}। (ग) $3114/300 = \qty{10.4}{km}$ प्रति मिलियन वर्ष,
\qty{1.0}{cm} प्रति वर्ष; $4926/300 = \qty{16.4}{km}$ प्रति मिलियन वर्ष, \qty{1.6}{cm} प्रति वर्ष।
\end{solution}

\begin{solution}{exo:b2:phanerozoic-earth:2}
(क) 72, 120.1, 50.5 और \qty{135.4}{Myr}; माध्य $(444 - 66)/4 = \qty{94.5}{Myr}$। (ख)
पर्मियन-अंत (साइबेरियाई, \qty{1.9}{Myr} का अंतर), ट्राइऐसिक-अंत (CAMP, \qty{0.4}{Myr}) और
क्रिटेशियस-अंत (डेक्कन, \qty{0.5}{Myr})। (ग) नहीं: अंतराल 50 से
\qty{135}{Myr} तक हैं; तिथियों के किसी भी समुच्चय का माध्य होता है।
\end{solution}

\begin{solution}{exo:b2:phanerozoic-earth:3}
(क) $\qty{91}{m}/\qty{0.036}{mm.yr^{-1}} = \num{2.5e6}$ वर्ष। (ख) $3.7/0.036 = 103$ गुना। (ग)
किसी भी तट से दूर, स्वच्छ जल में, नदियों की थोड़ी ही कीचड़ तल तक पहुँचती थी; शैवाल-कवचों की
वर्षा ही मुख्य अवसाद थी।
\end{solution}

\begin{solution}{exo:b2:phanerozoic-earth:4}
(क) नदी-बलुआ पत्थर से खड़िया तक निक्षेप ऊपर की ओर गहरे होते जाते हैं: एक ट्रांसग्रेशन।
खड़िया के ऊपर की समुद्र-तट की बालू एक रिग्रेशन दर्ज करती है। (ख) नहीं: वह सापेक्ष समुद्र-स्तर की एक वृद्धि दर्ज करता है,
जिसे स्थान का अवधसान, या अवसाद की आपूर्ति में कमी, भी उत्पन्न कर सकती थी। (ग)
वही ट्रांसग्रेशन, उसी समय, दूसरे महाद्वीपों पर खोजना, या \omterm{def:b2:basins-compared:backstripping}{बैकस्ट्रिपिंग} से
अवधसान हटाना।
\end{solution}

\begin{solution}{exo:b2:phanerozoic-earth:5}
(क) $233/97 = \qty{2.4}{m}$ प्रति मिलियन वर्ष, \qty{2.4}{\micro m} प्रति वर्ष। (ख) $3.7/0.0024 =
1540$: आज की वृद्धि कोई 1500 गुना तेज़ है। (ग) बैकस्ट्रिप किया वक्र: वेधन ने दिखाया है
कि चक्र-चार्ट गिरावटों का समय सही बताते हैं पर उनके परिमाण को दो या अधिक के गुणक से बढ़ाकर दिखाते
हैं; न्यू जर्सी का अपना पूर्वाग्रह है (गतिक स्थलाकृति), इसीलिए अनुमान $100 \pm
\qty{50}{m}$ है।
\end{solution}

\begin{solution}{exo:b2:phanerozoic-earth:6}
(क) $\num{e5} \times \num{e9} \times 2900 = \qty{2.9e17}{kg}$ बेसाल्ट, और $\num{2.9e17}
\times 0.0075 = \qty{2.2e15}{kg}$ कार्बन डाइऑक्साइड, $\num{2.2e15}/0.04401 = \qty{4.9e16}{mol}$।
(ख) $\num{278e-6} \times (44.01/28.97) \times \num{5.1e18} = \qty{2.15e15}{kg}$। (ग) एक स्पंद ने
लगभग उतना छोड़ा जितना 1750 में पूरी वायु में था। (घ) $\num{e17}/\num{2.15e15} = 46.5$: लगभग पचास।
\end{solution}

\begin{solution}{exo:b2:phanerozoic-earth:7}
(क) $2630 \times \frac{\pi}{6} \times (\num{e4})^3 = \qty{1.38e15}{kg}$। (ख) $\tfrac12 \times
\num{1.38e15} \times (\num{11.2e3})^2 = \qty{8.7e22}{J}$; \qty{20}{km.s^{-1}} पर,
\qty{2.76e23}{J}। (ग) \qty{47}{TW} का एक वर्ष $\num{47e12} \times \num{3.156e7} =
\qty{1.48e21}{J}$ है: 59 और 186 वर्ष।
\end{solution}

\begin{solution}{exo:b2:phanerozoic-earth:8}
(क) गहरा रंग, बारीक पटलन, बिलों का अभाव, बहुत कार्बनिक पदार्थ, पायराइट; प्रायः
कार्बन समस्थानिकों का विचलन और उष्णन के संकेत। (ख) टोआर्सियन घटना कारू--फ़ेरार के साथ, OAE~1a
ओंटोंग जावा के साथ। (ग) तिथियों में लाखों वर्षों तक की अनिश्चितताएँ होती हैं, और अनेक घटनाएँ तथा
प्रांत हैं, इसलिए संयोगवश मेल संभावित हैं; कारण-संबंध के लिए आवश्यक है कि प्रांत घटना से
ठीक पहले आरंभ हो और उन्हीं अवसादों में ज्वालामुखीयता का कोई अनुरेखक, जैसे पारा, मिले।
\end{solution}

\begin{solution}{exo:b2:phanerozoic-earth:9}
(क) $60.9 + 5.1 = \ang{66.0}$, $66.0 \times 111.2 = \qty{7339}{km}$, \qty{200}{Myr} में:
\qty{36.7}{km} प्रति मिलियन वर्ष, \qty{3.7}{cm} प्रति वर्ष। (ख) \qty{1.6}{cm} प्रति वर्ष की उस दर
के दुगुने से अधिक जो \qty{300}{Ma} से रही है। (ग) केवल अक्षांश का परिवर्तन ज्ञात है; देशांतर का
परिवर्तन, जो पुराचुंबकीय ढाँचे में मुक्त है, और कोई भी घुमावदार मार्ग दूरी में जुड़ते हैं।
\end{solution}

\begin{solution}{exo:b2:phanerozoic-earth:10}
(क) आयु $t$ के तल का अंश $p(t) = \frac2T(1 - t/T)$ है। तब $\int_0^T tp\,\mathrm{d}t =
\frac2T\bigl(\frac{T^2}2 - \frac{T^2}3\bigr) = \frac T3$ और $\int_0^T \sqrt t\,p\,\mathrm{d}t =
\frac2T\bigl(\frac23 T^{3/2} - \frac25 T^{3/2}\bigr) = \frac8{15}\sqrt T$। (ख) $T = 3 \times 64.2 =
\qty{192.6}{Myr}$; $2500 + \frac8{15} \times 350 \times 13.878 = \qty{5091}{m}$। (ग) $T/2 =
\qty{96.3}{Myr}$: \qty{4332}{m}; गिरावट \qty{759}{m} है:
\qty{2}{\percent} कम, \qty{774}{m} की तुलना में। (घ) समान मॉडल: $\frac23\sqrt\tau = \frac23\sqrt{2\bar t} = 0.943\sqrt{\bar t}$;
त्रिभुज: $\frac8{15}\sqrt{3\bar t} = 0.924\sqrt{\bar t}$। उत्तर आकार से कहीं अधिक माध्य आयु पर
निर्भर करता है।
\end{solution}

\begin{solution}{exo:b2:phanerozoic-earth:11}
(क) $\Delta s$ में केवल गुणक $1 - f^{-1/2}$ $f$ पर निर्भर है; $f = 2$ पर उसके मान,
$1 - 2^{-1/2} = 0.2929$, से भाग दीजिए। (ख) $f = (1 - 0.2929\,\Delta s/432)^{-2}$: \qty{50}{m} के लिए 1.071,
\qty{150}{m} के लिए 1.239; आज से 7 से \qty{24}{\percent} तेज़। (ग) $432 \times (1 -
0.8^{-1/2})/0.2929 = \qty{-174}{m}$: समुद्र \qty{174}{m} गिर जाता। (घ) आयु-वितरण
केवल तभी बदलता है जब पुराना तल क्षेपित होकर प्रतिस्थापित होता है, जिसमें $\tau$ की कोटि का समय लगता है,
लगभग \qty{100}{Myr}; समुद्र दसियों मिलियन वर्षों में उसका अनुसरण करता है।
\end{solution}

\begin{solution}{exo:b2:phanerozoic-earth:12}
(क) $2^{-7.4/2} = 0.077$: \qty{92}{\percent} गया। (ख) $t = 2\log_2(1/0.19) = \qty{4.8}{Myr}$।
(ग) $4.8/0.060 = 80$ गुना तेज़। (घ) एक लंबे चरण के भीतर, अकेला पृष्ठभूमि विलोपन अधिकांश
प्रजातियों को हटा देता है, इसलिए एक छोटा स्पंद एक बड़े योग में जुड़ता है और उससे अलग करना कठिन होता है;
अभिलेख की अपूर्णता भी अंतिम उपस्थितियों को पीछे की ओर फैला देती है।
\end{solution}

\begin{solution}{pb:b2:phanerozoic-earth:1}
\textbf{1.} 2500, 4250, 5304 और \qty{6000}{m}।
\textbf{2.} $\tau = 2 \times 64.2 = \qty{128.4}{Myr}$; $2500 + 350\sqrt{128.4} = \qty{6466}{m}$।
\textbf{3.} $\num{293e6}/(65000 \times 128.4) = \qty{35.1}{km}$ प्रति मिलियन वर्ष,
\qty{35.1}{mm} प्रति वर्ष: एक धीमी दर।
\textbf{4.} $\frac1\tau\int_0^\tau(2500 + 350\sqrt t)\,\mathrm{d}t = 2500 + \frac23 \times
350\sqrt\tau = \qty{5144}{m}$।
\textbf{5.} कुछ दसियों मिलियन वर्षों से पुराने तल के लिए, जो चपटा हो जाता है, वर्गमूल नियम बहुत गहरा है;
वास्तविक आयुएँ समान रूप से नहीं फैलीं, युवा तल सबसे अधिक है; और
अवसाद, \omterm{def:b2:oceanic-islands:seamount}{समुद्री पर्वत} और पठार तल को ऊपर उठाते हैं।
\textbf{6.} $\num{293e6} \times 5.144 = \qty{1.51e9}{km^3}$, पूरे महासागर से \qty{13}{\percent}
अधिक: मॉडल बहुत गहरा है।
\textbf{7.} $v = \qty{17.6}{km}$ प्रति मिलियन वर्ष; हर ओर
$17.6 \times 128.4 = \qty{2260}{km}$।
\textbf{8.} $t = 1000/17.6 = \qty{56.8}{Myr}$; $2500 + 350\sqrt{56.8} = \qty{5138}{m}$।
\textbf{9.} एक पार्श्व पर, $x = vt$ के साथ: $\int_0^{v\tau}\bigl(d(\tau) - d(x/v)\bigr)\mathrm{d}x
= vc\int_0^\tau(\sqrt\tau - \sqrt t)\,\mathrm{d}t = \frac13 vc\,\tau^{3/2}$; दोनों के लिए उसका दुगुना।
$\frac23 \times 17.6 \times 0.350 \times 128.4^{3/2} = \frac23 \times 17.6 \times 0.350 \times
1454.9 = \qty{5975}{km^2}$।
\textbf{10.} $65000 \times 5975 = \qty{3.88e8}{km^3}$।
\textbf{11.} $\num{3.88e8}/\num{1.33e9} = 0.29$।
\textbf{12.} उसका तल औसतन युवा है, इसलिए अधिक गर्म, कम घना और ऊँचा; अक्ष के अनुदिश युवा
तल की पट्टी अधिक चौड़ी है।
\textbf{13.} $\tau = \qty{64.2}{Myr}$; $2500 + \frac23 \times 350\sqrt{64.2} = \qty{4370}{m}$।
\textbf{14.} $5144 - 4370 = \qty{774}{m}$।
\textbf{15.} $\num{293e6} \times 0.774 = \qty{2.27e8}{km^3}$।
\textbf{16.} $\num{2.27e8}/\num{362e6} = \qty{0.627}{km}$, \qty{627}{m}।
\textbf{17.} $\Delta s = 627 \times (3300 - 1025)/3300 = \qty{432}{m}$।
\textbf{18.} $f = \bigl(1 - 0.2929 \times 100/432\bigr)^{-2} = 1.15$: \qty{15}{\percent} तेज़।
\textbf{19.} $0.149 \times 148 = \qty{22.1}{\mega km^2}$।
\textbf{20.} $48.1 + 0.32 \times 7.7 = \qty{50.6}{\percent}$, \qty{74.9}{\mega km^2}।
\textbf{21.} उठते हुए वह अधिक क्षेत्र ढकता है, इसलिए उतना ही आयतन एक पतली परत बनाता है
(\cref{prop:b2:phanerozoic-earth:flooding})।
\textbf{22.} 100, 200, 300 और \qty{400}{m} पर डूबे क्षेत्र: 22.1, 41.7, 57.3 और
\qty{71.2}{\mega km^2}। चरण: $(362 + 11.0) \times 0.1 = \qty{37.3e6}{km^3}$, फिर 39.4, 41.2
और \qty{42.6e6}{km^3}; संचयी 37.3, 76.7, 117.9 और \qty{160.5e6}{km^3}। रखा जाने वाला आयतन
\qty{156.4e6}{km^3} है, 300 और \qty{400}{m} के बीच: $300 + (156.4 - 117.9)/42.6 \times 100
= \qty{390}{m}$।
\textbf{23.} आज की भूमि का $38.7 + 0.90 \times 9.4 = \qty{47.2}{\percent}$।
\textbf{24.} प्रसरण-दर दुगुनी करने से समुद्र लगभग \textbf{\qty{390}{m}} उठता और
आज की भूमि का लगभग आधा भाग डूब जाता।
\end{solution}
